\chapter{方法}\label{chap2:methods}

\section{实验方法}\label{chap2:methods_exp}
\subsection{实验对象和材料}\label{chap2:methods_exp_materials}
本研究使用三只雄性恒河猴（Macaca mulatta，代号分别为 W、U、T，实验期间体重分别维持在：6.3--8.0千克，3.5--4.6千克，5.8--6.7千克）。所有实验程序均通过伦理委员会审查批准。

实验在隔音、暗光（仅屏幕光源）环境中进行，室温控制在20--25$^\circ$C。猕猴舒适地坐于定制猴椅中，头部固定，双眼正对约60厘米处的一台23.6英寸计算机屏幕，屏幕中心与双眼保持水平。视觉刺激呈现及奖赏发放通过基于MATLAB平台的MonkeyLogic 
2软件控制。

眼动数据采用EyeLink 1000红外眼动仪采集，采样频率为1000赫兹。神经元活动使用16通道线性阵列电极（Plexon S-probe）采集，电极通过多通道微操纵器（Alpha Omega EPS）推进至目标深度。信号由AlphaLab SnR系统（AlphaOmega公司）记录，采样率为44千赫兹，滤波范围为300--6000赫兹。

记录时段为：猕猴W，2023年1月至2023年7月；猕猴U，2024年2月至2025年3月；猕猴T，2023年9月至2024年9月。

\subsection{行为范式}\label{chap2:methods_exp_paradigm}

本研究训练三只猕猴完成图~\ref{fig2_1:task_paradigm}所示的赚取–消费任务。该范式的核心设计在于代币可跨试次累积：猕猴在每个试次中选择赚取代币或消费已累积的代币以换取饮用水奖励。赚取的代币存入钱包，每个试次剩余的代币结转至后续试次使用。赚取选项具有不同的赚取速率，消费选项具有不同的奖励量，两类参数在试次间独立随机变化（图~\ref{fig2_2:option_matrix}）。猕猴的最优策略取决于当前代币持有量与两个选项的具体参数，从而构成随财富状态动态变化的决策环境。

\begin{figure}[H]
    \centering
    \includegraphics[width=\textwidth]{Figure1_1_Paradigm}
    \bicaption{\enspace 赚取–消费范式}{\enspace The earning-spending paradigm}
    \label{fig2_1:task_paradigm}
\end{figure}
\fignotetext{屏幕中央的$5\times5$圆点阵列作为“钱包”持续显示当前代币数，在试次间隔期间保持可见。每个试次随机呈现两个几何图形，分别对应赚取选项与消费选项。猕猴通过眼跳至目标图形并维持注视完成选择：注视赚取选项时代币逐枚添加至钱包，注视消费选项时代币逐枚扣除并同步发放饮用水奖励。图中展示三次眼动，每次眼动可维持较长时间，期间可赚取或消费多枚代币，图中仅展示每次眼动选择的起始阶段。猕猴可在试次结束前随时切换至另一选项。试次时长从指数分布中随机采样，猕猴无法预测。灰色虚线指示注视位置，灰色箭头表示钱包代币数的变化方向；两者仅用于图示说明，在真实实验过程中不可见。}{The $5\times5$ circle array ("wallet") at the center of the screen continuously displays the current token count and remains visible during inter-trial intervals. On each trial, two geometric shapes were presented at randomized positions, corresponding to an earning option and a spending option. The monkey selected an option by making a saccade to the target shape and maintaining fixation; fixating the earning option added tokens to the wallet one at a time, while fixating the spending option removed tokens one at a time and simultaneously delivered a water reward. Three saccades are depicted; only the onset of each fixation is shown, as each fixation could last considerably longer in practice, during which multiple tokens were earned or spent. The monkey could freely switch between options at any time before trial end. Trial duration was sampled from an exponential distribution and was therefore unpredictable to the monkey. Gray dashed lines indicate fixation positions and gray arrows indicate the direction of wallet token changes; these gray elements are for illustration only and were not visible during the actual task.}

\begin{figure}[H]
    \centering
    \includegraphics[width=0.4\textwidth]{Figure1_2_Options}
    \bicaption{\enspace 选项条件矩阵}{\enspace Option condition matrix}
    \label{fig2_2:option_matrix}
\end{figure}
\fignotetext{赚取选项（三行）与消费选项（四列）交叉组合构成12种实验条件。颜色梯度表示各条件下赚取与消费的客观价值之比（赚取速率/奖励量）：绿色表示消费价值相对更高的条件，红色表示赚取价值相对更高的条件。}{Three earning options (rows) and four spending options (columns) are crossed to yield 12 experimental conditions. The color gradient represents the ratio of objective earning to spending value (earning speed / reward size): green indicates conditions in which spending value is relatively higher, and red indicates conditions in which earning value is relatively higher.}

屏幕中央始终呈现一个$5\times5$的圆形阵列，代表猕猴的"钱包"，其中实心圆的数量表示当前持有的代币数，阵列边长约$4.3^\circ$。代币数上限为25，下限为0：当持有代币数已达25时，选择赚取不会继续增加代币；当代币数为0时，选择消费不会继续扣除代币，但选择行为本身仍被记录。钱包阵列随赚取或消费行为实时更新，在试次间隔期间保持可见而选项刺激消失，以此强调代币的跨试次累积属性。

每个试次呈现两个简单几何图形，分别对应赚取与消费选项，猕猴需通过训练习得图形与选项属性之间的关联。两个选项呈现在距屏幕中心$6^\circ$处，位于试次开始时刻猕猴注视点方向顺时针与逆时针各$90^\circ$的对称位置，赚取选项与消费选项分居哪侧在每个试次中随机决定。这一设计在保证选项空间位置不可预测的同时，确保两个选项与起始注视点的距离始终相等，从而避免引入潜在的眼动成本差异。每个试次中，赚取速率与消费奖励量从各自参数集合中等概率随机抽取。实验共设置三种赚取速率：每250毫秒、200毫秒或100毫秒添加一个代币，对应速率分别为4、5、10代币/秒。消费选项共设置四种奖励量，以电磁阀开放时间作为饮用水体积的线性代理指标，电磁阀流速经标定约为0.2微升/毫秒；猕猴W和U的四种奖励量分别为0、20、35、55毫秒/代币，猕猴T的四种奖励量分别为15、25、35、70毫秒/代币。

猕猴需要注视目标并维持规定时长来完成选择。对于消费选项，每连续注视200毫秒即扣除一个代币（固定速率5代币/秒），同时发放该选项对应的饮用水奖励；对于赚取选项，代币按该选项的赚取速率逐次添加至钱包。猕猴可在试次内自由切换选项，切换通过眼跳至另一选项并维持注视完成。试次时长从指数分布中随机采样，均值为3000毫秒，截断范围为1000–5000毫秒，以减少动物对试次结束时刻的时间预期；试次间隔约1秒。每天实验开始时钱包代币数为0。

训练分为三个阶段，三只猕猴均按相同顺序完成：
\begin{enumerate}
    \item 注视训练。猕猴通过注视屏幕中央的注视点获得饮用水奖励。注视时长要求逐步延长，直至猕猴能够稳定维持注视1秒以上。进阶标准为连续3天注视成功率高于90\%。
    \item 眼跳训练。猕猴首先注视屏幕中央的注视点500毫秒；随后一个目标点在距屏幕中心$6^\circ$处随机方向呈现，猕猴需将注视迅速转移至目标点并维持注视1秒以获得奖励。进阶标准为连续3天成功率高于90\%。
    \item 赚取–消费任务。完成眼跳训练后，猕猴直接进入完整任务。训练初期将奖励量设定为较高值，以保证每日充足的水摄入，随后逐步递减至实验设定值。训练持续至行为表现稳定（模型估计的心理物理曲线在约一个月内无明显变化）后，开始采集神经元数据。
\end{enumerate}

三只猕猴均在约6个月内完成全部训练。

\subsection{手术与磁共振成像}\label{chap2:methods_exp_surgery}
行为训练前，每只猕猴接受头柱安装手术，用于实验期间的头部固定。

猕猴行为表现稳定后，接受第二次手术，在前额叶皮层上方植入丙烯酸材质的记录室。本文采用立体定位中常用的三个空间轴向：前后轴（anterior-posterior, AP）、内外轴（medial-lateral, ML）和上下轴（superior-inferior, SI）。记录室内径沿AP方向为25毫米，沿ML方向为20毫米，沿SI方向为10毫米。猕猴W和T的记录室置于左半球，猕猴U的记录室置于右半球。记录室内的颅骨被移除，暴露硬脑膜，并保持其结构完整，作为后续电生理记录的通道。

手术均在无菌条件下进行。猕猴术前经肌肉注射盐酸氯胺酮诱导麻醉（5–15毫克/千克），术中以1.5\%–2\%异氟烷维持麻醉。手术全程由专人持续监测猕猴的心率、血压、体温及呼气末二氧化碳等基本生命体征。

记录室植入后，对猕猴头部进行核磁扫描（Siemens MAGNETOM Prisma 3T），以确认记录室覆盖的脑区范围。扫描时将浓度为1摩尔/升四水合氯化锰母液稀释66倍后注入网格板盲孔作为造影剂。网格板盲孔孔径1毫米、间距2毫米的盲孔，其方向与电极穿透方向一致。记录位点的脑区归属综合两类信息进行推断：一是植入记录室后核磁影像中造影剂与大脑结构的空间对应关系，二是电极穿透过程中的电生理特征（包括放电密度及阻抗变化）。据此推定，本实验中ACC的记录范围对应BA 24区，OFC对应BA 11区，DLPFC对应BA 46区，各脑区的前后坐标范围标注于图~\ref{fig2:recording_sites}。共记录到679个ACC神经元（猕猴W：247；猕猴T：432；猕猴U无ACC记录）、1056个DLPFC神经元（猕猴W：127；猕猴U：491；猕猴T：438）以及903个OFC神经元（猕猴W：111；猕猴U：446；猕猴T：346）。

\begin{figure}[H]
    \centering
    \includegraphics[width=\textwidth]{Figure2_RecordingSites}
    \bicaption{\enspace 记录位点}{\enspace Recording sites}
    \label{fig2:recording_sites}
\end{figure}
\fignotetext{ACC（红色）、DLPFC（绿色）和OFC（蓝色）的记录位点分别叠加标注在三只猕猴的冠状面结构磁共振图像上：猕猴W（左）、猕猴U（中）、猕猴T（右）。前后坐标以耳棒零点为参考，标注于各截面左侧。本实验使用线性阵列电极记录，每个矩形标记代表一个记录位点，矩形沿电极轴方向的长度约为2毫米，略大于实际记录范围（约1.5毫米），以计入记录过程中可能发生的电极漂移与定位误差；重复记录的位点仅绘制一次。箭头标注主要脑沟的解剖名称。}{Recording sites for ACC (red), DLPFC (green), and OFC (blue) were overlaid on coronal structural MRI images for each monkey: Monkey W (left), Monkey U (middle), and Monkey T (right). Anterior–posterior (AP) coordinates relative to ear bar zero are indicated to the left of each section. Linear array electrodes were used; each rectangle represents one recording site, with the length along the electrode axis set to approximately 2~mm — slightly larger than the actual recording span ($\sim$1.5~mm) — to account for electrode drift and localization uncertainty. Sites recorded across multiple sessions are shown only once. Arrows indicate major sulci with their anatomical names.}

\section{行为数据分析}\label{chap2:methods_BHV}
所有记录了电生理数据的实验日（session）的行为数据均被纳入分析。为排除连续选择中惯性效应的干扰，所有分析均限定于每个试次序列中的首次选择；若首次选择未能成功执行（即未完成赚取或消费动作），则剔除该试次。此外，当财富水平达到钱包上限（25）或者下限（0）时，猕猴并无真正意义上的选择，对应试次亦不纳入后续分析。

\subsection{非参数主观权重估计}\label{chap2:methods_BHV_nonparametric}
为了在不预设特定函数形式的前提下量化财富水平对选项主观价值的调节作用，我们在全部24个财富水平（累积代币数量），分别估计了7种选项（包括3种赚取速度选项和4种消费水量选项）的主观权重。

在给定财富水平（$w$）、赚取速度（$e$）和消费水量（$s$）的条件下，选择“赚取”的概率 $P_{Earn}$ 可由以下函数建模：

\begin{equation}
    P_{Earn}(w,e,s) = \frac{1}{1 + e^{-\left(V^{Non}_E(w) - V^{Non}_S(w)\right)}}
\end{equation}

其中，$V^{Non}_E(w)$和$V^{Non}_S(w)$分别表示在财富水平$w$下，赚取选项与消费选项的估计权重。这两个权重反映了在当前财富状态下，各选项对选择“赚取”概率的相对贡献；上标“$^{Non}$”标识非参数（non-parametric）模型。该模型共包含$24 \times 7 = 168$个自由参数。

损失函数为交叉熵损失（cross-entropy loss）。对于第$i$次观测，假设真实选择为$c_i \in \{0, 1\}$（1 代表选择赚取，0 代表选择消费），模型预测的赚取概率为$p_i$，则数据拟合的损失函数为：

\begin{equation}\label{eq:cross_entropy}
    L_{data} = -\sum_{i} \left[c_i \ln p_i + (1-c_i) \ln (1-p_i)\right]
\end{equation}

为了防止模型过拟合，并引入“相邻财富水平下的选项主观权重不应剧烈跳变”的先验假设，我们在总损失函数中加入了跨财富水平的平滑惩罚项 $L_{smooth}$：

\begin{equation}
    L_{smooth} = \lambda \sum_{j=1}^{7} \sum_{w=1}^{23} (V_{j,w} - V_{j,w+1})^2
\end{equation}

其中，$j$ 遍历 7 种不同的选项权重，$w$ 遍历相邻的财富水平。$\lambda = 10$ 为经验选定的平滑惩罚强度。总损失函数为 $L = L_{data} + L_{smooth}$。

模型优化采用 MATLAB 的 \texttt{fminunc} 函数，基于拟牛顿算法（Quasi-Newton method）进行求解。为保障数值稳定性，预测概率被截断至 $[10^{-10}, 1 - 10^{-10}]$ 的区间。参数向量被初始化为微小的随机值以打破对称性，优化的终止条件设定为：最大迭代次数 2000 次、函数容差（Function tolerance）与步长容差（Step tolerance）均设为 $10^{-8}$。

\subsection{效用模型}\label{chap2:methods_BHV_utility}
非参数权重估计能够刻画财富调制效应的基本特征，但参数数量较多（$7$ 种选项形状 $\times 24$ 个财富水平）且权重本身缺乏可解释性。为揭示上述调制效应背后的计算机制，我们构建了一个基于边际效用理论的参数化模型，将核心假设浓缩至少量具有明确经济学含义的参数中。

模型假定效用函数为幂函数形式$u(w) \propto w^\alpha$，其中$w$表示当前财富（累积代币数），$\alpha$为边际效用曲率指数。选择幂函数形式是因为其简洁且不预设额外假设：$\alpha < 1$对应凹函数即边际效用递减，$\alpha > 1$ 对应凸函数即边际效用递增。而如果使用对数形式\citep{bernoulli1954exposition}，则无法兼容所有边际效用模式。在初始模型设定中，我们允许 $\alpha$ 在 $[0, +\infty)$ 范围内自由变化；但拟合结果显示$\alpha$均小于1，故后续拟合将$\alpha$约束在 $[0,1]$以提高数值稳定性。自由$\alpha$模型自然地嵌套了固定$\alpha$模型（$\alpha = 1$）作为特例，两者的比较可直接检验边际效用递减假设对行为解释力的贡献。

为定义选项的主观价值，首先需确定奖励量的度量。我们不直接使用电磁阀开放时间的物理数值，而是将奖励量归一化：以四种奖励量中第二小的值为单位$1$，其余奖励量按比例缩放。这一做法的动机在于假设猕猴经过训练后关注奖励量的相对差异而非绝对物理量；归一化后所有奖励量表达为无量纲的相对值$R$，省去了水量单位的细节，同时由于缩放差异会被交换率参数$ER$完全吸收，也不影响模型拟合。代币与归一化奖励之间的“汇率”$ER$（与每单位$R$等值的代币数）作为自由参数拟合。

在此基础上，赚取和消费选项的主观价值分别定义为：
\begin{align}
V_{E} &= \frac{ES}{ER} \cdot \alpha w^{\alpha - 1}\\
V_{S} &= R \cdot SS - \frac{SS}{ER} \cdot \alpha w^{\alpha - 1}.
\end{align}
其中$ES$ 为赚取速度（earning speed，单位：代币/秒）；$SS = 5$ 为消费速度（spending speed，单位：代币/秒）。

$V_{E}$ 的表达式表明，赚取一个额外代币的主观价值等于该代币的边际效用（$\alpha w^{\alpha - 1}$）乘以单位时间内的赚取速率（$ES/ER$，即每秒赚取的“奖励当量”）。$V_{S}$ 由两项构成：第一项 $R \cdot SS$ 为消费的即时奖励收益（每秒获得的归一化奖励量）；第二项为消费的财富成本。当 $R = 0$ 时，第一项为零，$V_{S}$ 变为负值，这准确反映了零奖励选项不被偏好的行为结果。当固定 $\alpha = 1$ 时，边际效用不随财富水平变化，对应的主观价值记为 $V_{E}^{\text{fixed}}$ 与 $V_{S}^{\text{fixed}}$。

综合决策变量为$\Delta V = V_{E} - V_{S} - b$，其中 $b$ 为常数偏移项。需注意，$b$ 与 $V_E$ 和 $V_S$ 中不依赖于财富和选项参数的常数成分在数学上不可区分；因此估计的主观价值与“真实”的主观价值之间可能相差一个常数，但对于仅关心相对水平的分析没有影响。选择概率由sigmoid函数给出：
\begin{equation}
P_{\text{earn}} = \frac{1}{1 + e^{-\frac{\Delta V}{\gamma}}},
\end{equation}
其中 $\gamma > 0$ 为选择敏感性，控制决策的随机程度：$\gamma$ 越小，选择越趋于确定；$\gamma$ 越大，选择越随机。

该拟合使用类似公式\eqref{eq:cross_entropy}的交叉熵损失函数估计参数 $\{\alpha, \, ER, \, b, \, \gamma\}$。优化采用MATLAB中的fmincon函数（Sequential Quadratic Programming, SQP算法），并配合多起点优化（MultiStart，10个随机起始点）以降低陷入局部最优的风险。固定 $\alpha$ 模型的$\alpha$固定为1，其余参数自由估计。

\subsection{留一交叉验证}\label{chap2:methods_BHV_LOO}
单个实验日内的试次在时间上相邻且共享同一行为状态；若仅在同一实验内划分训练集与测试集，可能高估模型的泛化表现。为严格检验自由$\alpha$模型是否优于固定$\alpha$模型，我们采用留一实验交叉验证（leave-one-session-out, LOSO），将整个实验日作为留存单元以避免信息泄漏。若自由$\alpha$模型在此条件下仍表现出预测优势，则该优势更可能反映真实的信号而非过拟合。

在每一折（fold）中，基于其余所有实验日的合并数据（训练集）估计模型参数，并在被留存的实验日（测试集）上计算模型表现。自由$\alpha$模型中的$\alpha$被约束在$[0,1]$区间。模型表现以测试集上的平均交叉熵量化：

\begin{equation}
    L_{test} = -\frac{1}{N_{\text{test}}} \sum_{i=1}^{N_{\text{test}}} \bigl[ c_i \ln \hat{p}_i + (1 - c_i) \ln(1 - \hat{p}_i) \bigr]
\end{equation}

其中$N_{\text{test}}$为测试集的总试次数，$c_i$为真实选择，$\hat{p}_i$由相应模型预测。模型能力的差异定义为$\Delta L = L^{fixed}_{test} - L^{free}_{test}$，正值表明自由$\alpha$模型的预测表现优于固定$\alpha$模型。

\section{神经数据分析}\label{chap2:methods_neural}
\subsection{预处理}\label{chap2:methods_neural_preprocess}
波形分类由Kilosort4算法\citep{pachitariu2024spike}自动完成。为评估自动分选的可靠性，我们随机抽取了部分神经元，将Kilosort4的分选结果与在线分选结果进行比对，二者在活跃时段、波形和放电率上高度一致，因此分类结果未再经过人工筛选。为排除电极漂移或细胞损伤造成的静默期，每个单元只保留其持续活跃的时段：活跃时段定义为在60秒时间窗内平均放电率（Firing rate，FR）高于1赫兹的时间区间。

原始数据首先对齐至每个试次中选项呈现的时刻（即试次开始时刻），随后以$10$毫秒的时间窗对数据进行分段，计算每个单元在每个时间窗内的放电率（firing rate，FR）。

对于涉及时间序列的分析，例如示例神经元的刺激时间直方图（peristimulus time histogram，PSTH）及信号编码的时间动态，放电率信号需经平滑处理。平滑采用因果高斯核$\sigma = 100$毫秒，半窗宽设为 $300$毫秒，核函数在时间零点截断以确保任一时刻的平滑放电率仅依赖于该时刻之前的神经活动。对于在特定时间窗内取平均后再进行分析的情形（如大多数单神经元选择性分析），放电率不再做额外平滑。

以下分析主要关注四个时间段的神经元活动：
\begin{itemize}
    \item 选项呈现：对齐到每次试次间隔结束后两个选项呈现在屏幕上的时刻，标志试次的开始；
    \item 钱包更新期（预期）：对齐到猕猴维持注视、钱包发生更新的时刻，回归变量为钱包更新后的财富水平，捕捉神经元对即将到来的财富状态的预期编码；
    \item 钱包更新期（反馈）：对齐到猕猴维持注视、钱包发生更新的时刻，回归变量为钱包更新前的财富水平，捕捉神经元对刚刚发生的选择结果的反馈编码；此阶段仅在图~\ref{fig4_6:update}相关分析中使用；
    \item 选择切换：对齐到猕猴发生选择切换的时刻。切换定义为：前后两次选择不同，且均成功触发代币的赚取或消费（在钱包处于上下限边界时，代币数量不变但选择行为仍被记录），从注视离开上一选项至到达下一选项的间隔不超过150毫秒。
\end{itemize}

\subsection{选择编码}\label{chap2:methods_neural_choice}
为量化单个神经元对动作域（左/右）和物品域（赚取/消费）选择信息的编码强度，对每个神经元计算选项呈现后0--500毫秒内的平均放电率。为消除不同神经元之间放电率的差异、使回归系数在神经元间可比较，对每个神经元的平均放电率进行z-score标准化。将标准化后的放电率作为因变量，对两个域的选择变量进行多元线性回归：
\begin{equation}
    FR \sim \beta_0 + \beta_1 \cdot Choice_{ES} + \beta_2 \cdot Choice_{LR},
\end{equation}
其中$Choice_{ES}$取$+1$（赚取）和$-1$（消费），$Choice_{LR}$取$+1$（左侧视野）和$-1$（右侧视野）。回归给出对应系数的$p$值，$p_{ES}$与$p_{LR}$分别量化该神经元对物品域和动作域的选择编码强度；$p_{ES} < p_{LR}$表示该神经元对物品域的选择编码强于对动作域的选择编码。

在群体水平，为考察神经元整体对两个域的编码偏向性，提取每个神经元回归系数的绝对值$|\beta_{ES}|$与$|\beta_{LR}|$，以二者构建散点图。计算$|\beta_{LR}|$与$|\beta_{ES}|$之间的 Spearman 秩相关系数，以检验动作域与物品域的编码强度是否显著相关。每个神经元向反对角线（$y=-x$）投影，投影值正比于神经元在两个域上编码强度的差值（$d = \frac{|\beta_{ES}| - |\beta_{LR}|}{\sqrt{2}}$）。在散点图内叠加该$d$值的分布直方图，并对其符号进行双侧二项检验：若$d$群体分布显著高于随机水平$0.5$，则表明神经元群体的表征整体偏向物品域；反之则偏向动作域。

\subsection{群体神经元活动降维}\label{chap2:methods_neural_TDR}

本研究中，财富水平、选项价值与选择类型等变量之间存在一定相关性。为在群体水平上解混相互关联的任务变量信号，我们采用了目标降维方法（Targeted Dimensionality Reduction，TDR；Mante等人\citep{mante2013computation}）。虽然解混主成分分析（Demixed PCA，dPCA，Kobak等人\citep{kobak2016demixed}）也能针对某一变量提取正交成分，但无法保证不同变量对应的主成分彼此正交。而 TDR 不同，TDR 通过正交三角分解（QR decomposition，以下简称QR分解）构建正交编码轴实现变量解混，其关键特性在于 QR 分解的列排序：排列靠前的变量优先占据协方差空间，靠后的变量仅解释剩余残差中的方差。若在某一靠后变量上检测到显著的编码信号，即表明该变量的群体编码在控制了相关变量后依然存在，而非源于变量间共线性的混淆。

本研究在标准TDR流程的基础上进行了两处改动。

其一，对神经元放电率采用逐时间点z-score归一化，而非将所有时间点合并后统一归一化。这一选择源于任务结构的考量：选项呈现阶段（包括视觉反应和运动信号）与保持注视阶段的放电率基线存在差异。若采用全局归一化，两阶段的基线偏移将叠加于投影结果之上；一旦该偏移与目标变量轴不正交，跨时间的编码强度比较便会引入系统性偏差，难以区分真实信号与归一化伪迹。逐时间点z-score等价于假设下游读出神经元对每个输入神经元的连接权重（$\beta$）与偏置项（bias）在每个时间点精确适配该神经元当前的均值与方差，这是一个比全局固定参数更强的假设。然而，这并不一定要求突触权重在毫秒尺度上发生可塑性变化——网络动态本身的时变演化可能可以完成上述适配。作为验证，我们也尝试用全局z-score重复了部分分析，结果趋势与主分析高度一致，仅在跨时间解码时观察到编码强度的下降，表明主要结论对归一化方案的选择具有稳健性。

其二，采用交叉验证策略以避免传统TDR中的正偏差问题。标准TDR流程以同一批数据同时定义编码轴与计算投影，投影结果天然偏向于高估编码强度。为此，我们将数据随机二等分，以训练集定义编码轴，以测试集计算投影。

\subsubsection{回归分析}

参与各分析的神经元须满足以下筛选标准：总有效试次数不少于50次，且每个条件分组至少包含10个有效试次。各时间窗内的放电率经平滑处理后独立进行 z-score 标准化，任务变量亦进行 z-score 标准化，以使不同变量的拟合系数具有可比性。随后对每个神经元拟合线性回归模型，估计系数向量 $\boldsymbol{\beta}_n \in \mathbb{R}^{V}$；将各神经元的系数向量按行堆叠，形成群体系数矩阵 $\mathbf{B} \in \mathbb{R}^{N \times V}$。

图\ref{fig3_3:trajectory}和图\ref{fig3_4:rsquared}的分析采用以下回归模型：
\begin{equation}
\mathrm{FR} \sim \sum_i \beta_i V_i + \beta_W \mathrm{Wealth} + \beta_{LR} \mathrm{Choice}_{LR} + \beta_{ES} \mathrm{Choice}_{ES}
\end{equation}
其中 $\mathrm{FR}$ 为选项呈现后 0--500毫秒内的平均放电率，$V_i \in \{V_L, V_R,V_E, V_S\}$ 分别为动作域选项价值（左/右）和物品域选项价值（赚取/消费），$\mathrm{Wealth}$ 为选项呈现时的财富水平（累积代币数，连续变量），$\mathrm{Choice}_{LR}$ 和 $\mathrm{Choice}_{ES}$ 分别为猕猴在动作域和物品域的选择。由于 $V_L$/$V_R$ 与 $V_E$/$V_S$ 共线（$V_L + V_R = V_E + V_S$），该模型添加了 L2 正则化（$\lambda = 5$，经验选定）以稳定系数估计；其余所有模型均使用标准 OLS。图\ref{fig3_2:scatter}的分析表明动作域与物品域编码相对独立，故后续分析不再纳入动作域信息，并以不带下标的 $\mathrm{Choice}$ 指代物品域选择。

图\ref{fig4_4:crossdecoding}、图\ref{fig4_6:update}的分析采用以下回归模型：
\begin{equation}
\mathrm{FR} \sim \beta_0 + \beta_E V_E + \beta_S V_S + \beta_C \mathrm{Choice} + \beta_W \mathrm{Wealth}
\end{equation}
其中 $\mathrm{FR}$ 为钱包更新前100到钱包更新后200毫秒内的平均放电率。

图\ref{fig6_4:conditioned}的分析首先依据猕猴的行为输出（赚取/消费）将试次分组，然后用两组试次独立拟合回归模型：
\begin{equation}
\mathrm{FR} \sim \beta_0 + \beta_E V_E + \beta_S V_S + \beta_w \mathrm{Wealth}
\end{equation}
其中 $\mathrm{FR}$ 为选项呈现后 0--500毫秒内的平均放电率。

图\ref{fig6_5:latency}和图\ref{fig6_6:comparison}的分析采用以下回归模型：
\begin{equation}
\mathrm{FR} \sim \beta_0 + \beta_w \mathrm{Wealth} + \beta_C \mathrm{Choice} + \beta_E V_E + \beta_S V_S + \beta_{\mathrm{Chosen}} V_{\mathrm{Chosen}}
\end{equation}
其中 $\mathrm{FR}$ 为选项呈现后 0--500 毫秒内的平均放电率，$V_{\mathrm{Chosen}}$ 为被选择选项的价值（详细定义见\ref{chap2:methods_neural_Vchosen}节）。

\subsubsection{正交归一化与投影}

群体系数矩阵 $\mathbf{B}$ 通过 QR 分解进行正交归一化：
\begin{equation}
    \mathbf{B} = \mathbf{Q}\mathbf{R}
\end{equation}
其中 $\mathbf{Q} \in \mathbb{R}^{N \times V}$ 的各列构成任务相关群体轴的正交归一基。目标变量的编码轴对应 QR 分解中的最后列，该列方向由之前列所张成子空间的正交补唯一确定，而与前列变量的内部排序无关。因此，各分析中只需将目标变量固定置于末位，其余变量的顺序不影响目标编码轴的结果。

随后在每个时间点，将经 z-score 标准化放电率矩阵 $\mathbf{X}(t)$ 投影到该正交基上，得到群体的时间序列信号：
\begin{equation}
    \mathbf{P}(t) = \mathbf{Q}^\top \mathbf{X}(t)
\end{equation}

\subsubsection{交叉验证}

若用同一批试次同时拟合回归系数和计算投影，即便在零信号条件下也会引入正向偏差，这对微弱信号的影响尤为显著。为消除这一偏差，每个神经元的试次被随机划分为两部分：训练集用于拟合 $\mathbf{B}$ 并导出正交归一轴 $\mathbf{Q}$，测试集用于计算投影。z-score 标准化的均值和标准差由训练集计算，并同步应用于测试集放电率，防止信息泄漏。

此外，各脑区记录的神经元数量差异较大（例如 DLPFC 有 1056 个，而 ACC 仅有 679 个），若直接使用全部神经元，分析结果（如以$R^2$衡量的编码强度）的估计方差将随脑区而变化，难以区分真实的编码差异与统计效力差异。为保证不同脑区和不同个体间降维结果的可比性，在每次交叉验证中对每个脑区无放回地固定抽取相同数量的神经元（个体数据分析：100 个；聚合数据分析：300 个），并对每个被抽中的神经元独立进行新的 50/50 试次划分。

为评估各时间点的编码强度，按条件将测试集试次分组，每组有放回重采样相同数量的试次。无论分析涉及的条件组数有多少，均抽取合计 300 个试次。将各组投影值与对应条件标签拼接为单一向量，编码强度定义为拼接投影向量 $\mathbf{p}(t)$ 与标签向量 $\mathbf{l}$ 的平方皮尔逊相关系数：
\begin{equation}
R^2(t) = \left[\operatorname{corr}\!\left(\mathbf{p}(t),\,\mathbf{l}\right)\right]^2
\end{equation}
上述流程在 $N_\mathrm{CV}$ 次自助法迭代中重复（主分析：100 次；跨时间解码：10 次）。

\subsubsection{跨时间解码}

为评估编码变量在任务不同时段的时间稳定性，我们对财富（$\mathrm{Wealth}$）进行跨时间解码。时间窗格为50毫秒，不做平滑处理，每个时间窗内的放电率取窗口内所有时间点的均值。分析中涉及的时间覆盖选项呈现期（$[-200,500]$毫秒，14 个时间窗）和钱包更新期（$[-200,300]$毫秒，10 个时间窗），组成 $24 \times 24$ 的解码矩阵；在每个训练时间窗，按前述回归和 QR 正交化步骤拟合 TDR 轴，并将其应用于所有测试时间窗计算 $R^2$（按财富水平分5组：[1]，[2-3]，[4-7]，[8-15]，[16-24]，每组 60 个试次）。该流程在 $N_\mathrm{CV} = 10$ 次自助法迭代中重复，迭代次数较主分析减少，以控制计算开销。

解码矩阵的每个元素均与打乱对照（shuffle control，在每个神经元内跨试次随机打乱变量标签，并重复完整的 TDR 流程得到的结果）进行比较。采用配对单侧 Wilcoxon 符号秩检验比较真实 $R^2$ 与打乱 $R^2$ 在各次迭代间的差异。

\subsubsection{潜伏期分析}

为衡量编码信号的时间特性，从每个 $R^2(t)$ 轨迹中提取两个时间指标：信号起始潜伏期（onset latency，$t_{onset}$）和半峰时间（$t_{50}$）。

首先以刺激前基线区间（$[-200,0)$毫秒）的 $R^2$ 值构建噪声分布，取其第 95 百分位数作为显著编码的阈值。

信号起始潜伏期（onset latency，$t_{onset}$）定义为： $R^2(t)$ 首次出现的持续至少 5 个连续时间点（$\geq 50$毫秒）均超过显著阈值的时间段的起始时刻；若不存在满足上述条件的时间点，则此次迭代的$t_{onset}$记为无效并排除。

半峰时间（$t_{50}$）定义为：在 $[0,\,T_\mathrm{peak}]$ 内，$R^2(t)$ 首次达到半峰值的时刻。半峰值定义为：
\begin{equation}
    R_\mathrm{half} = \frac{R_\mathrm{max} + R_\mathrm{baseline}}{2}
\end{equation}
其中$R_\mathrm{max}$为峰值，$R_\mathrm{baseline}$为基线中位数；若 $R_\mathrm{max}$ 未显著高于阈值，则此次迭代的 $t_{50}$ 记为无效并排除。

$t_{onset}$ 和 $t_{50}$ 的分布分别采用成对双侧 Wilcoxon 符号秩检验在变量间进行比较。

\subsection{财富编码}\label{chap2:methods_neural_wealth}
\subsubsection{财富调谐曲线}\label{chap2:methods_neural_wealth_tuning}
为量化单个神经元对财富状态的编码，取选项呈现期（ $[-300,300]$毫秒）平均放电率（未经平滑处理）作为因变量对任务变量进行多元线性回归：
\begin{equation}
    \mathrm{FR} \sim \beta_0 + \beta_E V_E + \beta_S V_S + \beta_C \mathrm{Choice}_{ES} + \beta_W \tilde{W}_5
\end{equation}
其中 $\tilde{W}_5$ 为选项呈现时被试的累积代币数，作为具有五个水平的类别变量（[1]、[2--3]、[4--7]、[8--15]、[16--24]）参与回归。对回归模型进行方差分析，若 $\tilde{W}_5$ 项的 $F$ 检验达到显著（$\alpha = 0.05$），则该神经元被归类为财富编码神经元。

偏好财富水平定义为在五组比较中产生最小 $p$ 值的水平，每组比较将某一财富水平的所有试次与其余四组合并后的试次进行双样本 $t$ 检验。调谐方向由偏好水平的平均放电率 $\bar{r}_{\mathrm{pref}}$ 与其余水平的平均放电率 $\bar{r}_{\mathrm{rest}}$ 之差的正负号决定。为在群体层面比较不同神经元的调谐形状，负向调谐神经元（$\bar{r}_{\mathrm{pref}} - \bar{r}_{\mathrm{rest}} < 0$）的调谐曲线在平均前乘以 $-1$，随后将所有神经元的调谐曲线施归一化：
\begin{equation}
    TC_{\mathrm{norm}} = \frac{TC - \min(TC)}{\max(TC) - \min(TC)}
\end{equation}
归一化曲线按具有相同偏好财富水平的神经元分组平均，得到各组的群体调谐曲线。

为刻画重新中心化后群体调谐曲线的形状，提取每个神经元偏好财富水平相对位置 $\pm 1$ 与 $\pm 2$ 处的归一化放电率，并以 Wilcoxon 秩和检验比较两者的分布。在零假设下，若调谐曲线仅反映随机噪声，则所有非峰值位置的放电率应无系统性差异；因此，$\pm 1$ 与 $\pm 2$ 位置间的显著差异说明峰值之外存在真实的梯度衰减。

对每个神经元独立执行 10 次打乱控制：每次将财富组标签跨试次随机置换后重复全部分析流程，10次结果的平均显著比例用于估计零假设下偏好财富水平的分布。

\subsubsection{对数压缩检验}\label{chap2:methods_neural_wealth_log}
为检验财富状态编码是否符合Weber-Fechner对数压缩规律，对每个纳入分析的财富编码神经元，分别在线性财富坐标（$[1,2,4,8,16]$）和对数财富坐标（$[1,2,3,4,5]$）下拟合Skew-Normal分布：
\begin{equation}
    f(w;\,\omega,\alpha) = \frac{2}{\omega}\,\phi\!\left(\frac{w-\xi}{\omega}\right)\Phi\!\left(\alpha\cdot\frac{w-\xi}{\omega}\right)
\end{equation}
其中$\phi$为标准正态概率密度函数，$\Phi$为标准正态累积分布函数，$\omega>0$为尺度参数，$\alpha$为偏态参数（$\alpha=0$时退化为对称Gaussian），位置参数$\xi$固定为0（即峰值对齐于偏好财富水平）。以残差平方和为损失函数，采用网格初始化结合\texttt{fmincon}优化求解$(\omega,\alpha)$。拟合结束后基于全局四分位距过滤离群值（$|\alpha - \text{median}| > 2.5 \times \text{IQR}$），随后以Wilcoxon符号秩检验分别检验线性和对数坐标下$\alpha$在神经元群体的分布与0的差异。

仅偏好财富水平落于中间三组（组2、3、4）的神经元被纳入本分析。边界组（组1和组5）的神经元因其调谐曲线在高值侧或低值侧被截断，Skew-Normal拟合所得偏态参数将系统性地偏向截断方向，无法可靠反映真实的编码不对称性，故予以排除。

\subsection{财富对选项价值编码的调制}\label{chap2:methods_neural_offer}
\subsubsection{联合编码分析}\label{chap2:methods_neural_offer_joint}
为刻画选项价值与财富在单神经元水平的共同表征，分别取选项呈现期（$[0,500]$毫秒）以及钱包更新期（$[-300,300]$毫秒）内的平均放电率（未经平滑）作为因变量，对每个神经元拟合以下回归模型：
\begin{equation}
    \mathrm{FR} \sim \beta_0 + \beta_E V_E^{\mathrm{fixed}} + \beta_S V_S^{\mathrm{fixed}} + \beta_C \mathrm{Choice} + \beta_W \mathrm{Wealth}
\end{equation}
其中 $V_E^{\mathrm{fixed}}$ 与 $V_S^{\mathrm{fixed}}$ 为固定 $\alpha$ 模型拟合得到的赚取与消费价值。使用客观价值而非主观价值，旨在确保价值项与财富项正交。若某变量对应的回归系数在 $p < 0.05$ 水平显著，且系数为正（负），则该神经元被归类为该变量的正向（负向）编码神经元。

对于每对变量（$V_E^{\mathrm{fixed}}$–Wealth 与 $V_S^{\mathrm{fixed}}$–Wealth），只有对两个变量均具有显著编码的神经元才进入 $2 \times 2$ 列联表，行与列分别对应两个变量各自的编码方向（正向/负向）。对于赚取价值，边际效用递减意味着财富越高，赚取选项的边际价值越低，因此预测正向编码 $V_E^{\mathrm{fixed}}$ 的神经元倾向于负向编码财富，反之亦然，列联表应该呈反对角集中模式（$V_E^{\mathrm{fixed}}$+/Wealth$-$ 与 $V_E^{\mathrm{fixed}}$-/Wealth$+$）；采用左尾 Fisher 精确检验该模式是否显著高于随机期望。对于消费价值，财富越高时消费的机会成本越低，主观消费价值随之上升，因此预测正向编码 $V_S^{\mathrm{fixed}}$ 的神经元倾向于正向编码财富，列联表应该呈对角集中模式（$V_S^{\mathrm{fixed}}$+/Wealth$+$ 与 $V_S^{\mathrm{fixed}}$-/Wealth$-$）；采用右尾 Fisher 精确检验该模式。

\subsubsection{嵌套模型比较}\label{chap2:methods_neural_offer_time}

为量化神经元是否编码了财富调制的价值信号，我们比较了两个回归模型对神经放电率的解释力。首先对平滑后的放电率矩阵在非重叠的100毫秒时间窗内取均值，分别覆盖选项呈现期（$[-200,800]$毫秒）、以及钱包更新期（$[-500,500]$毫秒）。钱包更新期的分析包含两个版本：全试次版本使用所有有效试次；试次匹配版本将有效试次从中随机下采样，使其与同一神经元在选项呈现期的有效试次数一致，以确保两个时间段的$\Delta R^2$在相同样本量下可比。

基线模型（$M_1$）包含选项价值、选择方向，以及财富水平：
\begin{equation}
    \mathrm{FR} \sim \beta_0 + \beta_E^{\mathrm{fixed}} V_E^{\mathrm{fixed}} + \beta_S^{\mathrm{fixed}} V_S^{\mathrm{fixed}} + \beta_C\,\mathrm{Choice} + \beta_W \tilde{W}_{24}
\end{equation}
其中$V_E^{\mathrm{fixed}}$和$V_S^{\mathrm{fixed}}$是固定$\alpha$模型拟合得到的不依赖于财富水平的赚取和消费价值，$\beta_W \tilde{W}_{24} = \sum_{k=2}^{24} \delta_k\,\mathbf{1}[Wealth = k]$将代币值1至24编码为24水平类别变量，以23个二元指示变量实现（代币值2至24各一个，以代币$= 1$为参考水平）。

财富调制模型（$M_2$）在$M_1$的基础上增加自由$\alpha$模型拟合得到的赚取价值$V_E$作为额外回归项：
\begin{equation}
    \mathrm{FR} \sim \beta_0 + \beta_E^{\mathrm{fixed}} V_E^{\mathrm{fixed}} + \beta_S^{\mathrm{fixed}} V_S^{\mathrm{fixed}} + \beta_C\,\mathrm{Choice} + \beta_W \tilde{W}_{24} + \beta_E V_E
\end{equation}

两个模型均以调整后$R^2$量化拟合优度，以惩罚$M_2$中额外的自由参数。财富调制指数$\Delta R^2 = R^2_{\mathrm{adj}}(M_2) - R^2_{\mathrm{adj}}(M_1)$在每个时间窗对每个神经元独立计算。

群体水平$\Delta R^2$的显著性采用双侧Wilcoxon符号秩检验，将每个神经元的$\Delta R^2$与打乱$V_E$得到的$\Delta R^2_{\mathrm{null}}$作比较。

\subsection{价值差编码}\label{chap2:methods_neural_deltaV}

为检验单神经元是否编码了价值差 $\Delta V = V_E - V_S$（此处省略常数偏移项$b$），我们分析 $V_E$ 与 $V_S$ 在单神经元水平的联合编码模式：若 $\Delta V$ 是主导编码量，则正向编码 $V_E$ 的神经元应倾向于负向编码 $V_S$，反之亦然，在 $2\times 2$ 列联表中形成反对角集中模式。

分别取选项呈现期（$[0, 500]$ 毫秒）与钱包更新期（$[0, 150]$ 毫秒）内的平均放电率（未经平滑）作为因变量，对每个神经元拟合以下回归：
\begin{equation}
    \mathrm{FR} \sim \beta_0 + \beta_E V_E + \beta_S V_S + \beta_C \mathrm{Choice} + \beta_W \tilde{W}_{24}
\end{equation}
由于某些代币值可能无对应试次，哑变量矩阵在稀疏数据下存在秩亏风险，回归采用伪逆（Moore--Penrose pseudoinverse）求解以稳健处理这一情形。神经元的分类标准与列联表构建规则同\ref{chap2:methods_neural_offer_joint}节。

采用双侧 Fisher 精确检验评估 $V_E \times V_S$ 列联表是否偏离独立性（$\alpha = 0.05$）。

\subsection{被选价值编码}\label{chap2:methods_neural_Vchosen}
行为模型中的常数偏移项 $b$ 与价值项（ $V_E$，$V_S$）中的常数成分无法区分，将其并入为 $V_E$ 或者 $V_S$ 将带有随意性。因此，我们定义：
\begin{equation}
    V_{\mathrm{Chosen}} = (V_E - V_S - b) \times c
\end{equation}
其中 $c = +1$ 为赚取选择，$c = -1$ 为消费选择，$b$ 为自由 $\alpha$ 行为模型中估计的偏移参数。$V_{\mathrm{Chosen}}$ 是被选选项相对于未选择选项的价值优势，但是仍然简称为被选价值。

将选项呈现期（$[0, 500]$ 毫秒）的平均放电率（未经平滑）作为因变量拟合以下回归模型：
\begin{equation}
    \mathrm{FR} \sim \beta_0 + \beta_E V_E + \beta_S V_S + \beta_{\mathrm{Chosen}} V_{\mathrm{Chosen}} + \beta_C \mathrm{Choice} + \beta_W \tilde{W}_{24}
\end{equation}
若 $\beta_{\mathrm{Chosen}}$ 的回归系数在 $p < 0.05$ 水平显著，则该神经元被归类为 $V_{\mathrm{Chosen}}$ 编码神经元。

每个神经元执行10次随机打乱对照分析，每次分析中分别计算显著神经元比例，最后取10次结果的平均值作为零假设下假阳性率估计。